% Abstract. Update the numbers here whenever a main result changes (see ../AGENTS.MD).

Nitrogen dioxide (\ce{NO2}) is a toxic air pollutant, and cheap sensors that work at room
temperature are needed to monitor it. Graphene quantum dots (GQDs) are small flakes of graphene
whose electronic structure depends on their size, shape and edges, so the right GQD could make a
sensitive and selective sensor material. The number of possible GQD shapes is very large and every
candidate needs quantum-chemical calculations, which makes an efficient search necessary. This
thesis builds and tests an agent workflow in which a large language model (LLM) proposes GQDs and
code checks and scores them.

The LLM does not write SMILES strings. It writes the positions of hexagons on a honeycomb grid,
and code turns them into a molecule, checks that the structure is closed-shell and valid, and,
when it is not, returns the reason together with edits that would fix it. With the closed-shell
rule in the prompt and fix edits in the feedback, 99\% of the LLM's designs were valid after at
most two retries, compared with 52\% without these aids. We then defined an \ce{NO2} sensing score
from GFN2-xTB calculations. It multiplies the charge transferred to \ce{NO2} by two factors: one
that falls to zero when \ce{NO2} binds too strongly to leave again at room temperature, and one
that requires \ce{NO2} to bind more strongly than CO, \ce{H2O} and \ce{NH3}. On 50 GQDs from an
existing dataset, the best scores came from GQDs with an xTB gap of about 0.35 to 0.7~eV that
physisorb \ce{NO2}. A DFT check (\(\omega\)B97X-V/def2-SVP) on two small models showed that xTB
describes physisorption well, overbinds chemisorbed \ce{NO2} and wrongly binds \ce{O2}.

The score was then used in an evolutionary loop in which either the LLM or random edits act as the
mutation operator, with the same number of xTB calculations for both. In a pilot with one run per
method, the LLM found the best GQD (score 0.340, against 0.331 for random edits and 0.322 at the
start), but the difference was not statistically significant. In the main comparison, with three runs per method, random edits found the best GQD in two
of the three pairs of runs (best score \(0.351 \pm 0.008\), against \(0.344 \pm 0.004\) for the
LLM), while the LLM's children scored higher on average in all three (\(0.186 \pm 0.012\) against
\(0.163 \pm 0.013\)). The LLM therefore proposes better edits on average but does not find better
GQDs than random edits. A DFT benchmark of twelve GQDs, run on a GPU, then tested the score
itself. xTB ranks the GQDs well by how strongly they bind \ce{NO2}, water and ammonia (Spearman
\(\rho = 0.74\) to 0.95), but not by charge transfer: in DFT, physisorbed \ce{NO2} takes almost no
charge (0.01~\(e\) against 0.23~\(e\) in xTB), so the score of every physisorbing GQD falls to about
zero and the score as a whole does not survive (\(\rho = -0.16\)). The best GQDs of both mutators
owe their high scores to this xTB artefact, so the response term of the score must be replaced.
DFT also shows that water, not CO, is the main competing gas, and that \ce{O2} does not bind.
Every proposal, accepted or rejected, is logged with its outcome, so the runs also
produce training data for fine-tuning a smaller model in future work.
